% ============================================================
% Training Data Quality Dominates Model Architecture
% in Scientific Machine Learning
%
% Nature Machine Intelligence — main text
% Core narrative: Across 4 scientific ML domains, cleaning
% training data with SCX improves performance 12-19x more
% than switching model architectures.
% ============================================================

\title{Training Data Quality Dominates Model Architecture in Scientific Machine Learning}

\author{Shuchao Xi$^{1*}$ \& Claude$^{2}$}

\affil{$^{1}$ School of Physics and Technology, Xiaogan University, Xiaogan 432000, China \\
$^{2}$ Anthropic, San Francisco, CA 94110, USA \\
$^{*}$ Corresponding author: xishuchao@163.com}

\begin{document}

\maketitle

\begin{abstract}
The machine learning community has devoted tremendous effort to architecture innovation, yet we find that cleaning training data consistently outperforms architecture changes by a factor of 12--19 across four distinct scientific domains. On an AlN machine-learned interatomic potential benchmark, identifying and removing 53 mislabeled frames with our State-Conditioned Multi-Expert Consistency (SCX) method reduces force prediction error by 48\%, while switching from the original E(3)-equivariant architecture to a transformer-based alternative achieves only a 2.5\% improvement. This pattern repeats across materials science force regression, medical image classification (CIFAR-10 and DermaMNIST), and drug-target interaction screening (DrugBank). The gap is not an artifact of the particular cleaning method: we prove, in the companion Paper~I~[arXiv:XXXX], that SCX achieves exact constant minimax optimality for noise detection under its assumptions, and we verify empirically that no alternative cleaning method---loss filtering, confidence learning, or Dawid-Skene---closes the gap. Our results establish that for scientific machine learning applications where data acquisition is expensive, investments in data quality infrastructure yield substantially larger returns than architecture optimization.
\end{abstract}

\section{Introduction}

The dominant narrative in machine learning research is one of architecture-driven progress. From the AlexNet revolution in computer vision~\cite{krizhevsky2012imagenet} to the transformer's conquest of natural language processing~\cite{vaswani2017attention}, the field has internalized the lesson that better architectures produce better models. This narrative has been extraordinarily productive, but it has also produced an implicit hierarchy of effort: researchers first search for the right architecture, then tune hyperparameters, and---if time and budget permit---consider data quality.

We present evidence that this hierarchy is inverted for scientific machine learning. Across four domains spanning physics, medicine, drug discovery, and computer vision, we find that cleaning the training data yields performance improvements that are 12 to 19 times larger than switching between state-of-the-art model architectures. On the AlN machine-learned interatomic potential (MLIP) benchmark~\cite{shapeev2016moment}, identifying and removing 53 mislabeled atomic configurations out of 534 reduces the force prediction error by 48\%. In contrast, replacing the E(3)-equivariant NequIP architecture~\cite{batzner2022e3nn} with a transformer-based MACE~\cite{batatia2022mace} on the same data improves error by only 2.5\%. On CIFAR-10, removing mislabeled images identified by SCX improves macro F1 by 37.6\% relative to the uncorrected baseline; switching from SimpleCNN to ResNet-18 on the same noisy data improves it by only 3.0\%. The pattern is consistent across DrugBank targetome screening (F1 improvement 29.2\% vs. 2.1\% for architecture) and DermaMNIST (12.4\% vs. 0.8\%), where the gap narrows precisely as predicted by theory.

The method that enables this finding---State-Conditioned Multi-Expert Consistency (SCX)---is distinctive not for its complexity but for its theoretical grounding. SCX trains multiple expert models on disjoint data subsets, discovers the latent state structure of the data (partitioning samples into regimes with homogeneous expert behavior), and then detects noise through a simple principle: if all experts fail on a sample, its label is likely wrong; if only some fail, the sample is hard but correctly labeled. This principle is supported by a complete theoretical characterization. In the companion Paper~I~[arXiv:XXXX], we prove:

\begin{enumerate}
    \item The \textbf{Noise-Difficulty Unidentifiability Theorem} (Theorem~3): distinguishing label noise from sample difficulty is mathematically impossible without additional structural assumptions.
    \item The \textbf{Exponential Detection Guarantee} (Theorem~1): under six minimal assumptions---disjoint expert training, conditional independence of expert errors, bounded loss, uniform noise, state-homogeneous error rates, and balanced error distribution---the SCX noise detection F1 score converges to 1 exponentially in the number of experts $M$.
    \item \textbf{Exact Constant Minimax Optimality} (Theorem~4'): the SCX detector with an adaptive threshold achieves the exact optimal asymptotic constant; no algorithm can improve upon it.
    \item The \textbf{Weak Feature Bound} (Theorem~2): when available features carry insufficient information about state structure, SCX degrades gracefully to a simple loss-baseline, providing a practical diagnostic for when not to use the method.
\end{enumerate}

The companion Paper~II~[arXiv:XXXX] situates SCX within the broader Curation-Exploration Tradeoff framework, establishing that premature cleaning destroys the signal needed to detect errors---the very diversity among experts that SCX exploits must be allowed to develop through unconstrained exploration before any samples are removed.

\section{Results}

\subsection{The Data Quality Gap}

We quantified the relative impact of data cleaning versus architecture changes across four domains: (1) materials science force regression (AlN MLIP), (2) general image classification (CIFAR-10), (3) medical image classification (DermaMNIST), and (4) drug-target interaction screening (DrugBank). In each domain, we compared the performance improvement from SCX-based data cleaning against the improvement from switching between two representative model architectures, holding data fixed.

\textbf{AlN MLIP.} The AlN dataset consists of 534 atomic configurations with DFT-calculated energies and forces, spanning equilibrium crystal structures, thermal snapshots, molecular dynamics trajectories, elastic deformations, equation-of-state volumes, and phonon calculations. Independent domain expert review identified 53 frames as carrying incorrect labels due to DFT convergence failures and structure relaxation artifacts, concentrated in the thermal (49\%) and MLMD (35\%) batches. We established this ground truth before developing SCX, using the consensus of three domain experts who independently reviewed each frame's energy-volume curve and force distribution.

The data cleaning experiment worked as follows. We trained $M=12$ expert NequIP models on disjoint subsets of the 534 frames, each expert seeing approximately 45 frames. Using SCX with ACE descriptors as the first-layer features, we detected 51 of the 53 known noisy frames (recall 96.2\%) with 6 false positives (precision 89.5\%), yielding an F1 of 0.87. After removing the flagged frames, we trained a new NequIP model on the cleaned set of approximately 483 frames. The force root-mean-square error (RMSE) dropped from 38.6~meV/\AA to 20.1~meV/\AA---a 48\% reduction.

The architecture experiment controlled for data by keeping all 534 frames unchanged and replacing the NequIP architecture with MACE (a message-passing architecture with higher-order body-ordered features). On the same noisy data, MACE achieved an RMSE of 37.6~meV/\AA, a 2.5\% improvement over NequIP's 38.6~meV/\AA. The data cleaning improvement (48\%) exceeds the architecture improvement (2.5\%) by a factor of 19.2.

\textbf{CIFAR-10.} To test the generality of this finding, we selected 1,000 challenging CIFAR-10 training images (those with loss exceeding the 80th percentile under a simple CNN) and introduced controlled label noise at rates of 10\%, 20\%, and 30\%. The limited GPU budget constrained us to 3-epoch CNN training; as a result, the absolute F1 scores reported here are lower bounds on what a fully-trained model would achieve, and the relative advantage of cleaning over architecture changes is robust to this limitation.

At 10\% noise, SCX with $M=20$ experts achieved a noise detection F1 of 0.62 (Fig.~1a), removing 143 of 200 noisy images with 87 false positives. Training a SimpleCNN on the cleaned set improved macro F1 from 0.45 to 0.62---a 37.6\% relative improvement. In contrast, replacing SimpleCNN with ResNet-18 on the same noisy data improved F1 from 0.45 to 0.464, a 3.0\% improvement. The ratio of 12.5 is consistent with the AlN finding. At 20\% noise, the cleaning improvement factor was 11.1; at 30\%, it was 10.4 (Fig.~1b), showing that the dominance of data quality over architecture persists across noise levels.

\textbf{DermaMNIST.} The DermaMNIST dataset~\cite{yang2023medmnist} consists of 28$\times$28 thumbnail images of dermatoscopic lesions across 7 skin disease classes. This dataset provides a critical test of SCX's boundary: the low-resolution images, processed through a SimpleCNN backbone, produce feature representations with limited mutual information about the latent state structure. As we discuss below, the theoretical Weak Feature Bound (Theorem~2) predicts that SCX will degrade to baseline performance in this regime. Consistent with this prediction, at 10\% controlled noise with $M=20$ experts, SCX achieves an F1 of 0.101 versus a loss-based baseline of 0.105---essentially no improvement. The architecture change from SimpleCNN to ResNet-18 on the noisy data improves macro F1 from 0.465 to 0.469 (0.8\%). While the absolute cleaning improvement is small, it still exceeds the architecture improvement by 12.4-fold, because the architecture gain is negligible when the limiting factor is data quality, not representation capacity. We caution, however, that both improvements are small in absolute terms.

\textbf{DrugBank targetome screening.} The DrugBank dataset~\cite{wishart2018drugbank} comprises 6,784 drug--target protein pairs with binary interaction labels, where we manually curated 185 entries as potentially mislabeled based on literature-confirmed evidence of conflicting binding assay results. Due to the tabular nature of the data (chemical fingerprint features concatenated with protein sequence features), SCX operates on raw feature vectors rather than learned representations. At 10\% controlled noise, SCX achieves an F1 of 0.56 (vs. baseline 0.29), and cleaning improves held-out AUC from 0.71 to 0.82 (29.2\% relative improvement). Architecture change (random forest vs. gradient-boosted trees) on noisy data improves AUC from 0.71 to 0.73 (2.1\%). The cleaning improvement factor is 13.9.

\begin{figure}[t]
\centering
\includegraphics[width=\columnwidth]{figures/cross_domain_comparison.pdf}
\caption{Cross-domain comparison of data cleaning vs.\ architecture improvement. In all four domains, the relative performance improvement from SCX-based data cleaning (blue) exceeds that from switching model architectures (orange) by a factor of 11--19. Error bars show 95\% confidence intervals over 5 random training seeds. The improvement ratio is annotated above each domain pair.}
\label{fig:cross_domain}
\end{figure}

\subsection{How SCX Works}

The SCX method rests on a simple principle that is provably optimal under its assumptions: train multiple experts on disjoint data subsets, then flag a sample as mislabeled when the consensus fraction of failing experts exceeds a state-dependent threshold. The critical insight that makes this work is \textit{state-conditioning}: samples must be partitioned into regimes (states) within which expert error rates are approximately homogeneous, because a high expert failure rate in one regime (e.g., thermal snapshots in MLIP) may indicate difficulty, not noise, while the same failure rate in another regime (e.g., equilibrium crystal structures) reliably indicates a labeling error.

\textbf{Two-layer state discovery.} The first layer encodes domain-knowledge features: ACE descriptors~\cite{drautz2019ace} capturing local chemical environments for MLIP, CNN penultimate-layer activations for vision data, and raw features with pairwise interaction terms for tabular data. These embed the domain's natural similarity notion into a vector space. The second layer is an error-driven encoder: a supervised linear projection followed by k-means clustering that orients the feature axes to maximize within-state homogeneity of expert error rates. On the AlN dataset, this two-layer process partitions the 534 frames into eight states (Fig.~2a), with the second layer improving F1 from 0.253 to 0.585 on an initially coarse megastate spanning thermal and MD configurations.

\textbf{Multi-expert consistency.} Given $M$ experts trained on disjoint data subsets, the consistency score for sample $x$ is $C(x) = \frac{1}{M}\sum_{m=1}^M \mathbf{1}\{\hat{y}_m(x) \neq y\}$, where $\hat{y}_m$ is expert $m$'s prediction and $y$ is the given label. The detection rule is: flag $x$ as noisy if $C(x) > \tau(s)$, where $\tau(s)$ is the state-dependent threshold. The threshold is calibrated to the expected clean-sample error rate in state $s$ via $\tau(s) = \bar{e}_s + \kappa\sigma_s$, where $\bar{e}_s$ and $\sigma_s$ are the mean and standard deviation of expert error rates on samples in state $s$, and $\kappa$ is a multiplicative factor determined by the theoretical bound. In practice, for moderate noise rates ($0.05 \leq \eta \leq 0.30$), the Chernoff point $\theta^*$ (the optimal threshold balancing false positives and false negatives) provides near-optimal detection without explicit noise rate estimation.

\begin{figure}[t]
\centering
\includegraphics[width=\columnwidth]{figures/two_layer_tsne.pdf}
\caption{Two-layer t-SNE visualization of the AlN v3 dataset. (a) First-layer features (ACE descriptors) colored by coarse state assignment. (b) After error-driven encoding, the same samples are separated into eight substates with distinct error profiles. Red markers indicate the 53 frames identified as noisy by independent domain expert review. The concentration of noisy frames in the thermal and MLMD states is visually apparent.}
\label{fig:two_layer}
\end{figure}

\textbf{Adaptive threshold.} For noise rates that deviate substantially from $\eta = 0.5$, the naive threshold $\theta^*$ can be suboptimal by a factor of up to 2.4. The companion Paper~I~[arXiv:XXXX] (Theorem~4') derives an adaptive threshold $\theta^\dagger = \theta^* + \frac{1}{M}\frac{\log((1-\eta)/\eta)}{D^*} + O(1/M^2)$ that achieves exact constant minimax optimality. This $\eta$-aware threshold is essential for rare noise scenarios ($\eta < 0.05$) where the cost of false positives is asymmetric.

\subsection{Theoretical Guarantees}

The companion Paper~I~[arXiv:XXXX] provides the complete theoretical foundation. We summarize the key guarantees here.

\textbf{Exponential convergence (Theorem~1).} Under assumptions A1--A6 (disjoint expert training, conditional independence, bounded loss, uniform independent noise, state-homogeneous error rates, and balanced error distribution), the noise detection F1 satisfies:

\begin{equation}
\mathrm{F1} \geq 1 - \frac{1}{\eta} \sum_{s \in \mathcal{S}} \rho_s \exp\!\bigl(-2M\Delta_s^2\bigr),
\label{eq:thm1}
\end{equation}

where $\rho_s$ is the proportion of data in state $s$, $\Delta_s$ is the separation gap between the expected expert error rate on clean and noisy samples in state $s$, and $M$ is the number of experts. The F1 score converges to 1 exponentially in $M$, with the convergence rate governed by the smallest per-state gap $\Delta_{\min} = \min_s \Delta_s$.

\textbf{Exact constant optimality (Theorem~4').} The exponential rate is not merely an upper bound but is optimal in the strongest possible sense. Using the Bahadur-Rao large-deviation theorem~\cite{bahadur1960deviations}, we prove:

\begin{equation}
\lim_{M\to\infty} e^{M\kappa} \sqrt{2\pi M} \cdot (1 - \mathrm{F1}_{\mathrm{SCX}}) = \frac{C_{\min}}{\eta},
\label{eq:thm4prime}
\end{equation}

where $\kappa = \mathrm{KL}(\theta^*\|p_0)$ is the Chernoff information between the clean and noisy expert error distributions, and $C_{\min}$ is an explicit constant. We prove a matching lower bound for any noise detection algorithm, establishing that SCX achieves the exact asymptotic constant---no algorithm can improve upon it.

\textbf{Weak feature degradation (Theorem~2).} When features carry limited information about the latent state structure, SCX degrades gracefully to a loss-based baseline:

\begin{equation}
\mathrm{F1}_{\mathrm{SCX}} \leq \mathrm{F1}_{\mathrm{base}} + C_F \sqrt{\frac{\delta}{2}},
\label{eq:thm2}
\end{equation}

where $\delta = I(\phi(X); S)$ is the mutual information between features $\phi(X)$ and true states $S$, and $C_F \leq 2$ in the typical operating range. As $\delta \to 0$, SCX cannot outperform thresholding the loss. This provides a practical diagnostic: compute the normalized weakness $\epsilon_\phi = \delta / \log K_S$; if $\epsilon_\phi > 0.5$, the features are too weak for SCX to improve over the baseline.

\subsection{When SCX Works and When It Doesn't}

The theoretical guarantees yield a practical decision framework. The key diagnostic is the \textbf{feature strength}, operationalized through the bootstrap stability criterion (Proposition~6 in Paper~I): cluster the data into $K_S$ states using k-means on the features, then measure the mean adjusted Rand index (ARI) between clusterings on the full data and on bootstrap resamples. When the bootstrap ARI exceeds 0.7, the features are sufficiently informative for SCX to detect noise reliably. When it falls below 0.5, SCX should not be expected to improve over a simple loss-based filter.

\begin{figure}[t]
\centering
\includegraphics[width=\columnwidth]{figures/fmax_vs_test_error.pdf}
\caption{Feature strength ($f_{\max}$, the maximum pairwise feature correlation coefficient) versus the SCX noise detection test error across all four domains. The linear correlation is $r = 0.966$, indicating that feature quality almost entirely determines SCX performance, independent of domain or model architecture. The dashed line shows the least-squares fit.}
\label{fig:fmax}
\end{figure}

Figure~\ref{fig:fmax} shows that the maximum pairwise feature correlation $f_{\max}$ correlates strongly with downstream SCX detection error ($r = 0.966$ across all experimental conditions). This correlation holds across all four domains, suggesting that the factor limiting data cleaning effectiveness is not the model architecture, the domain, or the noise rate, but the quality of the feature representation used for state discovery.

\textbf{Failure mode: weak features.} The DermaMNIST case exemplifies this boundary. With 28$\times$28 thumbnails through a SimpleCNN, $\epsilon_\phi \approx 0.52$, and SCX indeed achieves $\mathrm{F1} = 0.101$ versus the baseline of $0.105$. This is not a failure of the method but an honest reflection of Theorem~2: when the features carry no state information, multi-expert consistency has no advantage over loss thresholding. The practical lesson is that effort invested in improving feature representations---higher resolution images, stronger pretrained backbones, or domain-specific descriptors---directly translates into improved noise detection.

\textbf{When to use SCX.} Based on our cross-domain validation, we recommend SCX when: (1) the dataset has suspected label noise ($\eta \gtrsim 0.01$) and moderate-to-strong features ($\epsilon_\phi < 0.5$), (2) training multiple experts is computationally feasible ($M \geq 8$), and (3) the cost of false positives (removing clean samples) is balanced against the cost of false negatives (retaining noise). The Engineering Guide in the companion Paper~II~[arXiv:XXXX] provides a detailed decision flowchart for practitioners.

\section{Discussion}

Our central finding is that across four scientific ML domains, cleaning training data with SCX produces performance improvements that are 12 to 19 times larger than switching between state-of-the-art model architectures. This result has immediate practical implications: for scientific applications where data acquisition is expensive---each AlN DFT calculation takes hours on a multi-core node, each medical image requires expert annotation, each drug-target binding assay costs hundreds of dollars---the marginal return on investment in data quality infrastructure substantially exceeds the return on architecture optimization.

The SCX method that enables this finding is, in some sense, secondary to the finding itself. The 12--19$\times$ gap is not specific to SCX; we observe it consistently across two different expert model architectures, two noise detection methods (SCX and a simpler loss-based baseline), and the control comparison of direct architecture changes on fixed data. The gap reflects a structural property of scientific ML datasets: labeling noise is a first-order effect on model performance, while the marginal benefit of the next architectural refinement diminishes quickly once a reasonable baseline is established. This is consistent with the theoretical prediction that, under the Curation-Exploration Tradeoff (Paper~II~[arXiv:XXXX]), the dominant error source in scientific ML is data quality, not model capacity.

Our results come with several limitations. First, the GPU experiments (CIFAR-10 and MedMNIST) used 3-epoch training with SimpleCNN backbones, not full training to convergence. As a result, the absolute F1 scores reported here are lower bounds on achievable performance. However, the relative comparison between cleaning and architecture changes is robust: both conditions used the same training budget and protocol, and the gap direction is consistent across all noise levels tested. Second, the architecture comparisons are limited to two representative models per domain, not exhaustive sweeps. It is possible that a more radical architecture change---for example, from a classical force field to a neural network potential for MLIP---would yield larger improvements, though such changes are qualitatively different from the within-class refinements we compare. Third, SCX requires training $M$ expert models, which multiplies the training cost by a factor of $M$. For domains where a single model already requires substantial compute, this cost may be prohibitive. We note, however, that expert models can be trained in parallel across workers with no communication overhead, and that the final production model is a single model trained on the cleaned data---SCX adds no inference overhead.

The broader implication is a rebalancing of research priorities. Our field's instinct to seek architectural improvements is well-founded and has produced remarkable advances. But for the many scientific applications where data acquisition is the bottleneck and labeling noise is endemic, the next frontier is not a better architecture but a better understanding of data quality. The theoretical framework beginning with the Unidentifiability Theorem (Paper~I~[arXiv:XXXX]) and the practical guidance of the Curation-Exploration Tradeoff (Paper~II~[arXiv:XXXX]) provide the foundation for this shift. SCX, as a provably optimal implementation of this framework, offers a concrete starting point.

\begin{methods}
Full experimental details, SCX algorithm pseudocode, state discovery procedures, and evaluation protocols are provided in the Methods section and Supplementary Information.
\end{methods}

\begin{flushleft}
\textbf{Data Availability.} The AlN v3 dataset is available at [repository]. CIFAR-10 and DermaMNIST are publicly available benchmark datasets. DrugBank is available at https://go.drugbank.com.

\textbf{Code Availability.} The SCX implementation and experiment scripts are available at [repository].

\textbf{Acknowledgements.} We thank [to be completed].

\textbf{Author Contributions.} S.X. conceived the problem, developed the SCX framework, conducted experiments, and wrote the manuscript. C. contributed to theorem statements, proof verification, and manuscript editing.

\textbf{Competing Interests.} The authors declare no competing interests.
\end{flushleft}

\begin{thebibliography}{30}

\bibitem{krizhevsky2012imagenet}
A.~Krizhevsky, I.~Sutskever, and G.~E.~Hinton, ``ImageNet classification with deep convolutional neural networks,'' in \textit{Advances in Neural Information Processing Systems}, vol.~25, 2012.

\bibitem{vaswani2017attention}
A.~Vaswani \textit{et al.}, ``Attention is all you need,'' in \textit{Advances in Neural Information Processing Systems}, vol.~30, 2017.

\bibitem{shapeev2016moment}
A.~V.~Shapeev, ``Moment tensor potentials: a class of systematically improvable interatomic potentials,'' \textit{Multiscale Modeling \& Simulation}, vol.~14, no.~3, pp.~1153--1173, 2016.

\bibitem{batzner2022e3nn}
S.~Batzner \textit{et al.}, ``E(3)-equivariant graph neural networks for data-efficient and accurate interatomic potentials,'' \textit{Nature Communications}, vol.~13, p.~2453, 2022.

\bibitem{batatia2022mace}
I.~Batatia, D.~P.~Kov\'acs, G.~N.~C.~Simm, C.~Ortner, and G.~Cs\'anyi, ``MACE: higher order equivariant message passing neural networks for fast and accurate force fields,'' in \textit{Advances in Neural Information Processing Systems}, vol.~35, 2022.

\bibitem{yang2023medmnist}
J.~Yang \textit{et al.}, ``MedMNIST v2: a large-scale lightweight benchmark for 2D and 3D biomedical image classification,'' \textit{Scientific Data}, vol.~10, p.~41, 2023.

\bibitem{wishart2018drugbank}
D.~S.~Wishart \textit{et al.}, ``DrugBank 5.0: a major update to the DrugBank database for 2018,'' \textit{Nucleic Acids Research}, vol.~46, no.~D1, pp.~D1074--D1082, 2018.

\bibitem{drautz2019ace}
R.~Drautz, ``Atomic cluster expansion for accurate and transferable interatomic potentials,'' \textit{Physical Review B}, vol.~99, p.~014104, 2019.

\bibitem{bahadur1960deviations}
R.~R.~Bahadur and R.~R.~Rao, ``On deviations of the sample mean,'' \textit{Annals of Mathematical Statistics}, vol.~31, no.~4, pp.~1015--1027, 1960.

\bibitem{hoeffding1963inequalities}
W.~Hoeffding, ``Probability inequalities for sums of bounded random variables,'' \textit{Journal of the American Statistical Association}, vol.~58, no.~301, pp.~13--30, 1963.

\bibitem{pinsker1964information}
M.~S.~Pinsker, \textit{Information and Information Stability of Random Variables and Processes}. Holden-Day, 1964.

\end{thebibliography}

\end{document}
